% !PRINTED 331
\weberpage{331}
\begin{center}
\S.\,89.

\gesperrt{Inflexionspunkte einer Curve dritter Ordnung.}
\end{center}

Es sei jetzt $f=0$ die Gleichung einer Curve dritter Ordnung
 ohne singulären Punkt. Wir legen zwei der Ecken des
Coordinatendreiecks $\xi_1, \xi_2$ in zwei von den neun Inflexionspunkten
und nehmen die entsprechenden Inflexionstangenten für die
Seiten $x_2, x_1$. Dann muss sich $f$, wenn $x_1=0$ oder $x_2=0$ gesetzt
wird, auf das Glied $h\,x_3^3$ reduciren, worin $h$ eine Constante ist,
und folglich müssen alle Glieder, die in der cubischen Form
einen von Null verschiedenen Coëfficienten haben, ausgenommen
$h\,x_3^3$, durch $x_1\,x_2$ theilbar sein. Demnach erhält $f$ die Form

\weq{1}{f\,(x_1,x_2,x_3) = x_1\,x_2\,(a_1\,x_1+a_2\,x_2+a_3\,x_3) + h\,x_3^3,}

worin $a_1, a_2, a_3$ Constanten sind. Daraus aber geht hervor, dass
die gerade Linie

\[
a_1\,x_1 + a_1\,x_2 + a_3\,x_3 = 0
\] % [sic] printed coefficient of x_2 is a_1; expected a_2

gleichfalls Inflexionstangente der Curve $f$ ist, und dass ihr Schnittpunkt
 mit der Linie $x_3=0$ ein dritter Inflexionspunkt ist.
Dieser Punkt ist von den beiden ersten $\xi_1, \xi_2$ verschieden, weil
weder $a_1$ noch $a_2=0$ sein kann, wenn die Curve keinen singulären
 Punkt hat. Wir wollen dieses wichtige Resultat in Form
eines Satzes aussprechen:

1. \gesperrt{Wenn man irgend zwei Inflexionspunkte einer
Curve dritter Ordnung ohne singulären Punkt
durch eine gerade Linie verbindet, so schneidet
diese gerade Linie die Curve in einem dritten
Inflexionspunkte.}

Da es neun Inflexionspunkte giebt, so gehen von jedem
dieser Punkte vier gerade Linien aus, deren jede noch zwei
weitere Inflexionspunkte enthält. Es giebt neun solcher Büschel
von vier Geraden, und weil jede dieser Geraden in drei Büscheln
vorkommt, so ist die Anzahl der Geraden $9\,.\,4 : 3 = 12$. Wir
können demnach den Satz 1. so ergänzen:

2. \gesperrt{Die neun Inflexionspunkte einer Curve dritter
Ordnung liegen zu je dreien auf zwölf geraden
Linien, und durch jeden Inflexionspunkt gehen
vier von diesen Geraden.}
